% Part: proof-theory
% Chapter: sequent-calculus
% Section: invertibility

\documentclass[../../../include/open-logic-section]{subfiles}

\begin{document}

\olfileid{pt}{seq}{inv}

\olsection{规则的可逆性}

\begin{defn}
规则~$R$\[
\AxiomC{$S_1$}
\AxiomC{$\dots$}
\AxiomC{$S_n$}
\RightLabel{$R$}
\TrinaryInfC{$S$}
\DisplayProof
\]在系统中是\emph{可逆的}，若每当 $\Proves S$ 便有 $\Proves[h_i] S_i$（$i = 1$, \dots,~$n$）。它在系统中是\emph{高度保持可逆的}（或称\emph{hp-可逆的}），若每当 $\Proves_n S$ 便有 $\Proves[n] S_i$（$i = 1$, \dots,~$n$）。
\end{defn}

\begin{lem}\ollabel{lem:G3c-invert}
对于 $\lnot$、$\land$、$\lor$ 和 $\lif$，\Log{G3c} 的规则是高度保持可逆的，即：
\begin{enumerate}
  \item \RightR{\lnot}: 若 $\Log{G3c} \Proves[n] \Gamma \Sequent
  \Delta, \lnot !A$，则 $\Log{G3c} \Proves[n] !A, \Gamma \Sequent
  \Delta$。
  \item \LeftR{\lnot}: 若 $\Log{G3c} \Proves[n] \Gamma, \lnot !A
  \Sequent \Delta$，则 $\Log{G3c} \Proves[n] \Gamma \Sequent \Delta,
  !A$。
  \item \RightR{\land}: 若 $\Log{G3c} \Proves[n] \Gamma \Sequent
  \Delta, !A \land !B$，则 $\Log{G3c} \Proves[n] \Gamma \Sequent
  \Delta, !A$ 且 $\Log{G3c} \Proves[n] \Gamma \Sequent
  \Delta, !B$。
  \item \LeftR{\land}: 若 $\Log{G3c} \Proves[n] \Gamma, !A \land !B
  \Sequent \Delta$，则 $\Log{G3c} \Proves[n] !A, !B, \Gamma \Sequent
  \Delta$。
  \item \RightR{\lor}: 若 $\Log{G3c} \Proves[n] \Gamma \Sequent
  \Delta, !A \lor !B$，则 $\Log{G3c} \Proves[n] \Gamma \Sequent
  \Delta, !A, !B$。
  \item \LeftR{\lor}: 若 $\Log{G3c} \Proves[n] !A \lor !B, \Gamma \Sequent
  \Delta$，则 $\Log{G3c} \Proves[n] !A, \Gamma \Sequent
  \Delta$ 且 $\Log{G3c} \Proves[n] !B, \Gamma \Sequent
  \Delta$。
  \item \RightR{\lif}: 若 $\Log{G3c} \Proves[n] \Gamma \Sequent
  \Delta, !A \lif !B$，则 $\Log{G3c} \Proves[n] !A, \Gamma \Sequent
  \Delta, !B$。
  \item \LeftR{\lif}: 若 $\Log{G3c} \Proves[n] !A \lif !B, \Gamma
  \Sequent \Delta$，则 $\Log{G3c} \Proves[n] \Gamma \Sequent \Delta,
  !A$ 且 $\Log{G3c} \Proves[n] !B, \Gamma \Sequent \Delta$。
\end{enumerate}
\end{lem}

\begin{proof}
我们需要证明，对于 \Log{G3c} 的任一规则~$R$，若 $R$ 的结论~$S$ 存在一个推导~$\pi$，则 $R$ 的各前提 $S_i$ 也存在推导~$\pi_i$，且满足 $\pheight{\pi_i}
\le \pheight{\pi}$。以规则 \LeftR{\land} 为例：假设存在推导~$\pi$，其末相继式形如 $!A \land !B, \Gamma \Sequent \Delta$。我们需证明存在一个推导 $\pi'$ 使得 $!A, !B, \Gamma \Sequent \Delta$ 且 $\pheight{\pi'} \le \pheight{\pi}$。我们对 $n = \pheight{\pi}$ 进行归纳。

若 $n = 0$，则 $\pi$ 仅由一条公理组成：$!A \land !B, \Gamma \Sequent \Delta$。在 \Log{G3c} 中，公理形如 $!C, \Gamma \Sequent \Delta, !C$，其中 $!C$ 是原子公式。换言之，$!C$ 不可能是 $!A \land !B$。若 $!A \land !B, \Gamma \Sequent \Delta$ 是公理，则某个原子公式 $!C$ 必定同时属于 $\Gamma$ 和~$\Delta$。那么 $!A, !B, \Gamma \Sequent \Delta$ 也是一个公理，于是我们找到了所需的推导~$\pi'$（其高度也为~$0$）。

若 $n>0$，则 $\pi$ 的最后一步是一个推理。我们要证明断言对~$n$ 成立，但可归纳假设它对任何高度 $<n$ 的推导成立，即使上下文不同。换句话说，对任意 $k<n$ 以及任意 $\Pi$ 和 $\Lambda$，若 $!A \land !B, \Pi \Sequent \Lambda$ 有一个高度为~$k$ 的推导，则 $!A, !B, \Pi \Sequent \Lambda$ 有一个高度 $\le k$ 的推导。

我们分两种情况考虑：要么左侧的 $!A \land !B$ 是主公式，要么不是。如果它是主公式，则最后一步推理是 \LeftR{\land}，且其直接子推导~$\pi_1$ 以 $!A, !B, \Gamma \Sequent \Delta$ 结束。此时结论成立，因为 $\pheight{\pi_1} < \pheight{\pi}$。

如果 $!A \land !B$ 不是主公式，则必有其他某个公式是主公式，而 $!A \land !B$ 属于最后一步推理的上下文。那么归纳假设可应用于最后一条规则~$R$ 的直接子推导，从而得到一个推导~$\pi_1'$（或两个推导~$\pi_1'$ 和~$\pi_2'$），它们的高度不大于~$\pi$ 的直接子推导的高度。$\pi_1'$ 和 $\pi_2'$ 的末相继式中，左侧上下文包含的是 $!A, !B$ 而非 $!A \land !B$。应用规则~$R$ 即可得到推导~$\pi'$。例如，假设 $\pi$ 结束于 \LeftR{\lif}：
\[
\AxiomC{}
\RightLabel{$\pi_1$}
\Deduce$!A \land !B, \Gamma' \fCenter \Delta, !C$
\AxiomC{}
\RightLabel{$\pi_2$}
\Deduce$!A \land !B, \Gamma', !D \fCenter \Delta$
\RightLabel{\LeftR{\lif}}
\BinaryInf$!A \land !B, \Gamma', !C \lif !D \fCenter \Delta$
\DisplayProof
\]
这里左侧的 $!C \lif !D$ 是主公式，且左侧上下文是 $!A \land !B, \Gamma'$（即 $\Gamma = \Gamma', !C \lif !D$）。归纳假设分别给出 $!A \land !B, \Gamma' \Sequent \Delta, !C$ 和 $!A \land !B, \Gamma',
!D \Sequent \Delta$ 的推导 $\pi_1'$ 和 $\pi_2'$。这两个相继式可再次作为 \LeftR{\lif} 规则的前提，从而得到~$\pi'$：
\[
\AxiomC{}
\RightLabel{$\pi_1'$}
\Deduce$!A, !B, \Gamma' \fCenter \Delta, !C$
\AxiomC{}
\RightLabel{$\pi_2'$}
\Deduce$!A, !B, \Gamma', !D \fCenter \Delta$
\RightLabel{\LeftR{\lif}}
\BinaryInf$!A, !B, \Gamma', !C \lif !D \fCenter \Delta$
\DisplayProof
\]
我们有
\begin{align*}
\pheight{\pi} & = \max(\pheight{\pi_1}, \pheight{\pi_2}) + 1\\
\pheight{\pi'} & = \max(\pheight{\pi_1'}, \pheight{\pi_2'}) + 1
\intertext{（根据定义），且}
\pheight{\pi_1'} & \le \pheight{\pi_1}\\
\pheight{\pi_2'} & \le \pheight{\pi_2}
\intertext{（根据归纳假设）。因此，}
\pheight{\pi'} & \le \pheight{\pi}.
\end{align*}
\end{proof}

\begin{prob}
给出 \olref{lem:G3c-invert} 中关于 \RightR{\lif} 的证明。
\end{prob}

\begin{lem}\ollabel{lem:invert-quant}
规则 \RightR{\lforall} 和 \LeftR{\lexists} 在如下意义下是高度保持可逆的：
\begin{enumerate}
  \item 若 $\Log{G3c} \Proves[n] \Gamma \Sequent \Delta,
  \lforall[x][!A(x)]$，则 $\Log{G3c} \Proves[n] \Gamma \Sequent
  \Delta, !A(c)$；且
  \item 若 $\Log{G3c} \Proves[n] \lexists[x][!A(x)], \Gamma \Sequent
  \Delta$，则 $\Log{G3c} \Proves[n] !A(c), \Gamma \Sequent \Delta$，
\end{enumerate}
其中 $c$ 是任意不在 $\Gamma$、$\Delta$ 以及（分别对应）$\lforall[x][!A(x)]$ 或 $\lexists[x][!A(x)]$ 中出现的常项符号。
\end{lem}

\begin{proof}
  我们证明 (1)，将 (2) 留作练习。证明与 \olref{lem:G3c-invert} 类似。在归纳步骤中，同样有两种情形。第一种情形，即 $\lforall[x][!A(x)]$ 是主公式的情形，处理很简单：最后一个 $\RightR{\lforall}$ 推理的前提具有所需形式 $\Gamma \Sequent \Delta, !A(a)$，且结束于它的直接子推导~$\pi_1$ 满足 $\pheight{\pi_1} < \pheight{\pi}$。此外，本征变元条件告诉我们，$a$ 不在 $\Gamma$、$\Delta$ 或 $\lforall[x][!A(x)]$ 中出现。由于我们可以假设 $\pi_1$ 是正则的，那么根据 \olref[qua]{lem:subst-regular}，$\Subst{\pi_1}{c}{a}$ 是 $\Gamma \Sequent \Delta, !A(c)$ 的一个相同高度的推导。

  在第二种情形中，$\lforall[x][!A(x)]$ 不是最后一个推理的主公式；主公式是其他某个公式。我们再次以一个情形为例进行说明。假设最后一个推理是~\RightR{\land}。$\Delta$ 中某个形如~$!B \land !C$ 的公式是主公式。推导~$\pi$ 结束于
  \[
\AxiomC{}
\RightLabel{$\pi_1$}
\Deduce$\Gamma \fCenter \Delta', \lforall[x][!A(x)], !B$
\AxiomC{}
\RightLabel{$\pi_2$}
\Deduce$\Gamma \fCenter \Delta', \lforall[x][!A(x)], !C$
\RightLabel{\RightR{\land}}
\BinaryInf$\Gamma \fCenter \Delta', \lforall[x][!A(x)], !B \land !C$
\DisplayProof
\]
其中 $\Delta = \Delta', !B \land !C$。设 $c$ 是一个不在 $\Delta$ 中出现的常项符号；那么显然它也不在 $\Delta', !B$ 或 $\Delta', !C$ 中出现。因此归纳假设可应用于 $\pi_1$ 和~$\pi_2$；我们得到满足 $\pheight{\pi_1'} \le \pheight{\pi_1}$ 和 $\pheight{\pi_2'} \le
\pheight{\pi_2}$ 的推导~$\pi_1'$ 和 $\pi_2'$。所需的 $\Gamma \Sequent \Delta, !A(c)$ 的推导~$\pi'$ 为：
\[
\AxiomC{}
\RightLabel{$\pi_1'$}
\Deduce$\Gamma \fCenter \Delta', !A(c), !B$
\AxiomC{}
\RightLabel{$\pi_2'$}
\Deduce$\Gamma \fCenter \Delta', !A(c), !C$
\RightLabel{\RightR{\land}}
\BinaryInf$\Gamma \fCenter \Delta', !A(c), !B \land !C$
\DisplayProof
\]
它满足 $\pheight{\pi'} = \max(\pheight{\pi_1'}, \pheight{\pi_2'})+1 \le \pheight{\pi}$。
\end{proof}

\Olref{lem:G3c-invert} 在切割消去定理的证明中起到重要作用。但它也可以用来证明以下命题：

\begin{prop}\ollabel{prop:G3c-cont-adm}
系统~\Log{G1c} 的收缩规则 \LeftR{\Contraction} 和 \RightR{\Contraction} 在系统~\Log{G3c} 中是高度保持可接纳的。
\end{prop}

\begin{proof}
我们同时给出 \RightR{\Contraction} 和 \LeftR{\Contraction} 的论证。假设 $\pi$ 是 \Log{G3c} 中 $\Gamma \Sequent \Delta, !A, !A$ 或 $!A, !A, \Gamma \Sequent \Delta$ 的一个推导。我们需要证明，相应地也存在 \Log{G3c}-推导~$\pi'$，使得 $\Gamma \Sequent \Delta, !A$ 或 $!A, \Gamma \Sequent \Delta$ 成立，且满足 $\pheight{\pi'} \le \pheight{\pi}$。我们通过对 $n = \pheight{\pi}$ 施加归纳来证明这一点。

如果 $n=0$，则 $\pi$ 就是一条公理。但如果 $\Gamma \Sequent \Delta, !A, !A$ 是 \Log{G3c} 的公理，那么 $\Gamma \Sequent \Delta, !A$ 也是公理：要么某个原子公式~$!C$ 同时是 $\Gamma$ 和~$\Delta$ 的元素，要么 $!A$ 是原子的且是~$\Gamma$ 的元素。当末尾相继式为 $!A, !A, \Gamma \Sequent \Delta$ 时，同样的推理也适用。

如果 $n>0$，则 $\pi$ 以某条规则~$R$ 结束。如果 $!A$ 不是最后这个推理的主公式，我们可以像 \olref{lem:G3c-invert} 的证明中那样，将归纳假设应用于~$\pi$ 的直接子推导，再对所得推导应用同样的规则~$R$，从而得到推导~$\pi'$。归纳假设是：如果存在高度为~$k<n$ 的、形如 $\Pi \Sequent \Lambda, !D, !D$ 或 $!D, !D, \Pi \Sequent \Lambda$ 的推导，那么相应地也存在高度~$\le k$ 的、形如 $\Pi \Sequent \Lambda, !D$ 或 $!D, \Pi \Sequent \Lambda$ 的推导。（注意，即使上下文或被收缩的公式不同于 $\Gamma$、$\Delta$ 或~$!A$，归纳假设也是适用的！）

例如，假设 $R$ 是 $\RightR{\land}$ 且 $\pi$~的结尾为
\[
\AxiomC{}
\RightLabel{$\pi_1$}
\Deduce$\Gamma \fCenter \Delta', !B, !A, !A$
\AxiomC{}
\RightLabel{$\pi_2$}
\Deduce$\Gamma \fCenter \Delta', !C, !A, !A$
\RightLabel{\RightR{\land}}
\BinaryInf$\Gamma \fCenter \Delta', !B \land !C, !A, !A$
\DisplayProof
\]
其中 $\Delta = \Delta', !B \land !C$。归纳假设分别给出 $\Gamma \fCenter \Delta', !B, !A$ 和 $\Gamma \fCenter \Delta', !C, !A$ 的推导~$\pi_1'$ 和~$\pi_2'$，且满足 $\pheight{\pi_1'} \le  \pheight{\pi_1}$ 和 $\pheight{\pi_2'} \le
\pheight{\pi_2}$。则推导~$\pi'$ 为：
\[
\AxiomC{}
\RightLabel{$\pi_1'$}
\Deduce$\Gamma \fCenter \Delta', !B, !A$
\AxiomC{}
\RightLabel{$\pi_2'$}
\Deduce$\Gamma \fCenter \Delta', !C, !A$
\RightLabel{\RightR{\land}}
\BinaryInf$\Gamma \fCenter \Delta', !B \land !C, !A$
\DisplayProof
\]

如果 $!A$ 是最后这个推理的主公式，我们先将反转引理（\olref{lem:G3c-invert}）应用于~$\pi$ 的直接子推导，然后应用归纳假设，最后再次应用规则~$R$。例如，假设 $\pi$ 证明了 $\Gamma \Sequent \Delta, !A, !A$，其中 $!A \ident (!B \land !C)$ 且是最后这个推理的主公式。那么 $\pi$ 形如下式：
\[
\AxiomC{}
\RightLabel{$\pi_1$}
\Deduce$\Gamma \fCenter \Delta, !B \land !C, !B$
\AxiomC{}
\RightLabel{$\pi_2$}
\Deduce$\Gamma \fCenter \Delta, !B \land !C, !C$
\RightLabel{\RightR{\land}}
\BinaryInf$\Gamma \fCenter \Delta, !B \land !C, !B \land !C$
\DisplayProof
\]
对 $\pi_1$ 应用 \olref{lem:G3c-invert}，得到 $\Gamma \Sequent \Delta, !B, !B$ 的推导~$\pi_1'$。（它同时也产生 $\Gamma \Sequent \Delta, !C, !B$ 的推导，但我们不需要。）类似地，对~$\pi_2$ 应用 \olref{lem:G3c-invert}，我们得到 $\Gamma \Sequent \Delta, !C, !C$ 的推导~$\pi_2'$。由于 \RightR{\land} 的反转是保持高度的，故 $\pheight{\pi_1'} \le  \pheight{\pi_1} < \pheight{\pi}$ 和 $\pheight{\pi_2'} \le  \pheight{\pi_2} < \pheight{\pi}$。于是，归纳假设可应用于 $\pi_1'$ 和 $\pi_2'$：分别存在 $\Gamma \Sequent \Delta, !B$ 和 $\Gamma \Sequent \Delta, !C$ 的推导~$\pi_1''$ 和~$\pi_2''$，且满足 $\pheight{\pi_1''}
\le \pheight{\pi_1'} \le \pheight{\pi_1}$ 和 $\pheight{\pi_2''} \le
\pheight{\pi_2'} \le \pheight{\pi_2}$。则推导~$\pi'$ 即为：
\[
\AxiomC{}
\RightLabel{$\pi_1''$}
\Deduce$\Gamma \fCenter \Delta, !B$
\AxiomC{}
\RightLabel{$\pi_2''$}
\Deduce$\Gamma \fCenter \Delta, !C$
\RightLabel{\RightR{\land}}
\BinaryInf$\Gamma \fCenter \Delta, !B \land !C$
\DisplayProof
\]
其中 $\pheight{\pi'} \le \pheight{\pi}$。

当公式为主公式时，对于量词规则的情形，我们必须特别小心处理。

假设 $!A \ident \lforall[x][!B(x)]$ 在右侧是主公式。那么 $\pi$ 以如下形式结束：
\[
\AxiomC{}
\RightLabel{$\pi_1$}
\Deduce$\Gamma \fCenter \Delta, \lforall[x][!B(x)], !B(a)$
\RightLabel{\RightR{\lexists}}
\UnaryInf$\Gamma \fCenter \Delta, \lforall[x][!B(x)], \lforall[x][!B(x)]$
\DisplayProof
\]
根据 \olref{lem:invert-quant}，对于任何不在 $\Gamma \fCenter \Delta,
\lforall[x][!B(x)], !B(a)$ 中出现的~$c$，存在高度 $\le \pheight{\pi_1}$ 的推导~$\pi_1'$，其末端相继式为 $\Gamma \fCenter \Delta, !B(c),
!B(a)$。我们可以假定 $\pi_1'$ 是正则的；因此根据 \olref[qua]{lem:subst-regular}，推导 $\Subst{\pi_1'}{a}{c}$ 也是 $\Gamma \fCenter \Delta, !B(a),
!B(a)$ 的推导，其高度同样 $\le \pheight{\pi_1}$。现在应用归纳假设，得到 $\Gamma \fCenter \Delta, !B(a)$ 的推导~$\pi_1''$。再应用 \RightR{\lforall} 规则得到~$\pi'$：
\[
\AxiomC{}
\RightLabel{$\pi_1''$}
\Deduce$\Gamma \fCenter \Delta, !B(a)$
\RightLabel{\RightR{\lforall}}
\UnaryInf$\Gamma \fCenter \Delta, \lforall[x][!B(x)]$
\DisplayProof
\]
其中 $\pheight{\pi'} \le \pheight{\pi}$。

现在假设 $!A \ident \lexists[x][!B(x)]$ 且在右侧是主公式。那么 $\pi$ 形如下式：
\[
\AxiomC{}
\RightLabel{$\pi_1$}
\Deduce$\Gamma \fCenter \Delta, \lexists[x][!B(x)], \lexists[x][!B(x)], !B(t)$
\RightLabel{\RightR{\lexists}}
\UnaryInf$\Gamma \fCenter \Delta, \lexists[x][!B(x)], \lexists[x][!B(x)]$
\DisplayProof
\]
归纳假设应用于~$\pi_1$，存在 $\Gamma \Sequent \Delta, \lexists[x][!B(x)],
!B$ 的推导~$\pi_1'$。（注意这里不需要反转引理！）则推导~$\pi'$ 即为：
\[
\AxiomC{}
\RightLabel{$\pi_1'$}
\Deduce$\Gamma \fCenter \Delta,  \lexists[x][!B(x)], !B$
\RightLabel{\RightR{\lexists}}
\UnaryInf$\Gamma \fCenter \Delta, \lexists[x][!B(x)]$
\DisplayProof
\]
其中 $\pheight{\pi'} \le \pheight{\pi}$。
\end{proof}

\begin{prob}
概述 \olref{prop:G3c-cont-adm} 中关于 \LeftR{\Contraction} 的证明。描述 $!A \ident (!B \land !C)$ 和 $!A \ident (!B \lif !C)$ 的情形。
\end{prob}

\begin{prob}
概述 \olref{prop:G3c-cont-adm} 中关于 \LeftR{\Contraction} 在其余量词情形下的证明，即当 $!A \ident
\lforall[x][!B(x)]$ 和 $!A \ident \lexists[x][!B(x)]$ 时。
\end{prob}

\end{document}